% Compiled and edited by Brian Bi
\documentclass[10pt]{extarticle}
\setlength{\parindent}{0.0in}
\usepackage{xeCJK}
\usepackage{fontspec, xunicode, xltxtra}
\usepackage{amsmath}
\usepackage{amssymb}
\usepackage{amsfonts}
\usepackage{pifont}
\usepackage{hyperref}
\usepackage[a5paper,twoside=true,top=15mm,bottom=10mm,left=10mm,right=10mm]{geometry}
\usepackage{fancyhdr}
\pagestyle{fancy}
\fancyhead{}
\fancyfoot{}
\fancyhead[LE,RO]{\thepage}
\fancyhead[LO]{\leftmark}
\fancyhead[RE]{\rightmark}

\usepackage{listings}
\makeatletter
\lst@CCPutMacro\lst@ProcessOther {"2D}{\lst@ttfamily{-{}}{-{}}}
\@empty\z@\@empty
\makeatother
\usepackage{color}
\newfontfamily\mono{Consolas}
\newfontfamily\timesnew{Times New Roman}
\lstset {
  basicstyle = \small\mono,
  columns = fixed,
  extendedchars=false,
  %language = C++,
  tabsize = 4,
  frame = single,
  breaklines = true,
  breakindent = 1.1em,
  numbers=left,
  stringstyle=\mono,
  numberstyle=\footnotesize\mono,
  basewidth={0.5em, 0.4em},
}

\newcommand\Cross{%
 \makebox[1.0em][c]{%
 \makebox[0pt][c]{\raisebox{0em}{\ding{56}}}}~~}%

\newcommand\Warning{%
 \makebox[1.0em][c]{%
 \makebox[0pt][c]{\raisebox{0.1em}{\tiny!}}%
 \makebox[0pt][c]{$\bigtriangleup$}}~~}%

\newcommand\Tick{%
 \makebox[1.0em][c]{%
 \makebox[0pt][c]{\raisebox{0em}{$\checkmark$}}}~~}%

\newcommand\Require[1]{\ref{#1} \nameref{#1} }

\newcommand*\newevenpage{%
  \clearpage
  \ifodd\value{page}\hbox{}\newpage\fi
}

\newcommand{\BookNo}{1}
\newcommand{\BookTitle}{Graph Theory}
\begin{document}
\begin{titlepage}
\begin{center}
% Upper part of the page
\includegraphics[width=0.15\textwidth]{../shared/njulogo.pdf}\\[1cm]
\textsc{\huge Nanjing University}\\[3cm]
\timesnew{\huge ACM-ICPC Codebook \BookNo}\\[0.5cm]
% Title
%\HRule \\[0.4cm]
{ \Huge \bfseries \BookTitle}\\[0.4cm]
%\HRule \\[1.5cm]


\vfill

% Bottom of the page
{\large \today}

\end{center}

\end{titlepage}

\setmainfont{Times New Roman}
\setCJKmainfont{SimSun}
\setlength{\parskip}{0.0in}
\tableofcontents
\setlength{\parskip}{0.1in}
\newevenpage
% New codebook
\section{Shortest Paths}

\subsection{Single-source shortest paths}
\subsubsection{Dijkstra} \label{dijkstra}
Dijkstra's algorithm with binary heap. \par
\Cross Can't be performed on graphs with negative weights.\par
\textbf{Usage:} \\[0.1cm]
\begin{tabular}{p{2.5cm} p{9cm}}
  \lstinline|V| & Number of vertices. \textbf{PLEASE set this value before use!} \\
  \lstinline|add_edge(e)| & Add edge $e$ to the graph. \\
  \lstinline|dijkstra(src)| & Calculate SSSP from $src$.  \\
  \lstinline|d[x]| & distance to $x$ \\
  \lstinline|p[x]| & last edge to $x$ in SSSP  \\
\end{tabular} \par
\textbf{Time complexity:} $O(E \log V)$ \par
\lstinputlisting[language=c++]{shortest-paths/dijkstra.cpp}

\subsubsection{SPFA}
Shortest path faster algorithm. (Improved version of Bellman-Ford algorithm) \par
This code is used to replace \lstinline|void dijkstra(int src)|. \par
\Tick Can be performed on graphs with negative weights. \\
\Warning For some specially constructed graphs, this algorithm is very slow. \par
\textbf{Usage:} \\[0.1cm]
\begin{tabular}{p{2.5cm} p{9cm}}
  \lstinline|spfa(src)| & Calculate SSSP from $src$.  \\
\end{tabular} \par
\textbf{Requirement:} \\
\Require{dijkstra} \par
\textbf{Time complexity:} $O(kE)$, for most graphs, $k<2$ \par
\lstinputlisting[language=c++]{shortest-paths/spfa.cpp}

\subsection{All-pairs shortest paths (Floyd-Warshall)}
Floyd-Warshall algorithm. \par
\Tick Can be performed on graphs with negative weights. To detect negative cycle, one can inspect the diagonal, and the presence of a negative number indicates that the corresponding vertex lies on some negative cycle. \\
\Warning \textbf{Self-loops} and \textbf{multiple edges} must be specially judged. \\
\Warning If the weights of edges might exceed \lstinline|LLONG_MAX / 2|, the line \lstinline|(*)| should be added. \par
\textbf{Usage:} \\[0.1cm]
\begin{tabular}{p{2.5cm} p{9cm}}
  \lstinline|init()| & Initialize the distances of the edges from 0 to V.  \\
  \lstinline|floyd()| & Calculate APSP. \\
  \lstinline|d[i][j]| & distance from $i$ to $j$ \\
\end{tabular} \par
\textbf{Time complexity:} $O(V^3)$ \par
\lstinputlisting[language=c++]{shortest-paths/floyd.cpp}

\section{Spanning Tree}
\subsection{Minimum spanning tree}
\subsubsection{Kruskal's algorithm}
\textbf{Usage:} \\[0.1cm]
\begin{tabular}{p{2.5cm} p{9cm}}
  \lstinline|n, m| & The number of vertices and edges, resp. \\
  \lstinline|edges[]| & Edges of the graph, numbered from 0. \\
  \lstinline|kruskal()| & Run Kruskal's algorihtm. \\
\end{tabular} \par
\textbf{Time complexity:} $O(E \log E)$ \par
\lstinputlisting[language=c++]{spanning-tree/kruskal.cpp}

\subsubsection{Prim's algorithm, adjacency matrix representation}
Calculate minimum spanning tree. The result is represented as a tree rootes at \lstinline|src|. \par
\textbf{Usage:} \\[0.1cm]
\begin{tabular}{p{2.5cm} p{9cm}}
  \lstinline|adj[i][j]| & Adjacency matrix, indexed from 1. \\
  \lstinline|prim(src)| & Run Prim's algorihtm from \lstinline|src|. \\
\end{tabular} \par
\textbf{Time complexity:} $O(V^2)$ \par
\lstinputlisting[language=c++]{spanning-tree/prim.cpp}

\subsection{Minimum ratio spanning tree}
Minimize $\frac{\sum_{e \in ST} cost[e]}{\sum_{e \in ST} dist[e]}$ where $ST$ is a spanning tree. \par
\textbf{Usage:} \\[0.1cm]
First, build the edges of the graph as the structure shows; then, implement a usual MST algorithm; finally, call \lstinline|solve()| to get the answer. \par
\lstinputlisting[language=c++]{spanning-tree/ratio.cpp}

\subsection{Manhattan distance minimum spanning tree}
\textbf{Usgae:} \\[0.1cm]
\begin{tabular}{p{3cm} p{9cm}}
  \lstinline|add_point(x, y)| & Add point $(x, y)$. \\
  \lstinline|Manhattan_MST()| & Calculate Manhattan distance minimum spanning tree.
\end{tabular} \par
\textbf{Time complexity:} $O(n \log n)$, but constant factor may be large. \par
\lstinputlisting[language=c++]{spanning-tree/manhattan.cpp}

\section{Depth-first Search}
\subsection{Strongly connected components, condensation (Tarjan)}
Find strongly connected components and compute the component graph. \par
\Warning The component graph may contain \textbf{multiple edges}. \par
\textbf{Usage:} \\[0.1cm]
\begin{tabular}{p{2.5cm} p{9cm}}
  \lstinline|V| & number of vertices \\
  \lstinline|scc[i]| & the SCC that $i$ belongs to, numbered from 1. \\
  \lstinline|sccn| & number of SCCs \\
  \lstinline|find_scc()| & Find all SCCs. \\
  \lstinline|contract()| & Compute component graph. \\
\end{tabular} \par
\textbf{Time complexity:} $O(V + E)$ \par
\lstinputlisting[language=c++]{dfs/scc.cpp}

\subsection{Vertex biconnected components (Tarjan)}
Find all vertex biconnected components (blocks). \par
\textbf{Usage:} \\[0.1cm]
\begin{tabular}{p{2.5cm} p{9cm}}
  \lstinline|add_edge(u, v)| & Add edge $(u, v)$ to the graph. \textbf{Vertices are numbered from 0.}\\
  \lstinline|find_bcc(n)| & Find all BCCs in the graph containing $n$ vertices $[0, n)$.  \\
  \lstinline|bcc_cnt| & number of BCCs \\
  \lstinline|bccno[]| & the id of BCC a vertex belongs to.\textbf{The id is numbered from 0. If the vertex is a cut vertex, the value is undefined.} \\
  \lstinline|iscut[]| & whether a vertex is a cut vertex \\
  \lstinline|bcc[]| & the vertex set of a BCC \\
  \lstinline|bcce[]| & the edge set of a BCC \\
\end{tabular} \par
\textbf{Time complexity:} $O(V + E)$ \par
\lstinputlisting[language=c++]{dfs/bcc.cpp}

\section{Flow Network}
\subsection{Maximum flow, Dinic's algorithm}
\Tick Can be performed on networks with parallel and antiparallel edges.\par
\textbf{Usage:} \\[0.1cm]
\begin{tabular}{p{3.5cm} p{8cm}}
  \lstinline|add_edge(u, v, c)| & Add an edge from $u$ to $v$ with capacity $c$. \\
  \lstinline|max_flow(s, t)| & Compute maximum flow from $s$ to $t$. \\
\end{tabular} \par
\textbf{Time complexity:} For general graph, $O(V^2E)$; for network with unit capacities, \\ $O(\min\{V^{2/3}, E^{1/2}\}E)$; for bipartite network, $O(\sqrt{V}E)$. \par
\lstinputlisting[language=c++]{flow/dinic.cpp}

\subsection{MCMF, Ford-Fulkerson with SPFA}
\Tick Can be performed on networks with parallel and antiparallel edges.\par
\textbf{Usage:} \\[0.1cm]
\begin{tabular}{p{4.5cm} p{7cm}}
  \lstinline|add_edge(u, v, c)| & Add an edge from $u$ to $v$ with capacity $c$. \\
  \lstinline|min_cost(s, t, &cost)| & Compute MCMF from $s$ to $t$. Return the maximum flow, and set \lstinline|cost| as the minimum cost. \\
  \lstinline|min_cost(s, t, f, &cost)| & Compute minimum cost flow from $s$ to $t$ with capacity $f$. Return whether such flow exists, and set \lstinline|cost| as the minimum cost, if exists. \\
\end{tabular} \par
\textbf{Time complexity:} $O(|f|kE)$. \par
\lstinputlisting[language=c++]{flow/mcmf.cpp}

\section{Matching}
\subsection{Maximum cardinality bipartite matching, Hungarian}
\textbf{Usage:} \\[0.1cm]
\begin{tabular}{p{2.5cm} p{8cm}}
  \lstinline|init(nx, ny)| & Initialize the algorithm with two parts containing  \lstinline|nx|,  \lstinline|ny| vertices, respectively. The vertices are numbered from 0. \\
  \lstinline|add(a, b)| & Add an edge from $a$ to $b$. \\
  \lstinline|match()| & Compute and return maximum cardinality bipartite matching,. \\
  \lstinline|mx[], my[]| & The index of the other vertex if matched; otherwise $-1$. \\
\end{tabular} \par
\textbf{Time complexity:} $O(VE)$. \par
\lstinputlisting[language=c++]{matching/hungarian.cpp}
\end{document}
